\section{Introducción}

Los sistemas embebidos permiten integrar sensores, actuadores, interfaces de usuario y protocolos de comunicación para automatizar procesos físicos. En el Proyecto 1 de Electrónica Digital 2 se requiere diseñar e implementar una red de sensores formada por cuatro microcontroladores en total, utilizando tres ATmega328P y un ESP32. El sistema debe incorporar tres sensores distintos, de los cuales al menos uno debe utilizar comunicación I2C, además de una pantalla LCD, tres tipos de motores ---DC, Stepper y Servo--- y comunicación con la plataforma Adafruit IO mediante WiFi.

Para cumplir estos requerimientos se propone un \textit{Car Wash Inteligente}, representado mediante una maqueta que automatiza el recorrido de un automóvil de juguete a través de diferentes etapas de un lavado. El proceso comienza con la detección del vehículo en la entrada y continúa con su desplazamiento por las estaciones del sistema. Durante el recorrido se utilizan sensores para detectar la presencia del automóvil, monitorear el nivel de agua y determinar la llegada del vehículo al final del trayecto. De forma complementaria, distintos actuadores permiten controlar la barrera de entrada, el movimiento del automóvil y el accionamiento de los cepillos.

La arquitectura utiliza un ATmega328P como microcontrolador Maestro y dos ATmega328P como Slaves conectados mediante un bus I2C compartido. El Maestro también controla una pantalla LCD de 16x2 y se comunica con un sensor de distancia VL53L0X mediante I2C. Las mediciones y estados relevantes son transferidos al ESP32 mediante UART. El ESP32 implementa la conexión WiFi y el enlace con Adafruit IO, permitiendo publicar telemetría y, en modo manual, enviar órdenes de control hacia el Maestro.

El presente informe describe progresivamente el diseño, construcción e implementación del sistema. Las secciones de pruebas, resultados y análisis se completarán a partir de los datos obtenidos durante la implementación real, evitando incorporar resultados o mediciones no verificadas.
